\chapter{Project Manifests and Package Transformations}
\label{chap:manifests}

\epigraph{``Tell me what your software depends on, and I will tell you who you are: a reproducible wizard or an imperative gambler.''}{--- The Guix Hacker's Catechism}

\noindent
Typing \texttt{guix shell python python-numpy gcc-toolchain make git ...} every time you open a terminal quickly becomes tedious. Furthermore, passing command-line arguments does not capture custom package variants or programmatic dependencies.

To turn your project's development environment into a version-controlled, shareable artifact, GNU Guix provides \textbf{Manifests} and \textbf{Package Transformations}.

\section{Writing a \texttt{manifest.scm}}

A manifest is a small Guile Scheme file that returns a list of packages to be included in an environment.

There are two primary ways to write a manifest:

\subsection{1. The Specification Approach: \texttt{specifications->manifest}}

If you simply want a list of packages by name, \texttt{specifications->manifest} is the most concise:

\begin{lstlisting}[style=guixscheme]
;; manifest.scm - Simple specification manifest
(specifications->manifest
  '("python"
    "python-numpy"
    "python-pandas"
    "python-matplotlib"
    "jupyter"
    "gcc-toolchain"
    "make"))
\end{lstlisting}

\subsection{2. The Programmatic Approach: \texttt{packages->manifest}}

Because manifests are written in a full general-purpose programming language, you can construct package lists using Scheme logic, conditionals, and module imports:

\begin{lstlisting}[style=guixscheme]
;; manifest.scm - Programmatic manifest
(use-modules (gnu packages python)
             (gnu packages python-science)
             (gnu packages python-xyz)
             (gnu packages gcc)
             (gnu packages base)
             (guix packages))

(packages->manifest
  (list python
        python-scipy
        python-pytorch
        gcc-toolchain
        gnu-make))
\end{lstlisting}

\section{Automatic Manifest Discovery}

When you run \guixcmd{guix shell} in a directory containing a \texttt{manifest.scm} file without specifying package arguments, Guix automatically detects and loads the manifest:

\begin{lstlisting}[style=guixcli]
$ ls
manifest.scm  src/  Makefile

$ guix shell
[env]$ # All packages from manifest.scm are automatically active!
\end{lstlisting}

If you want to run your automated tests in a hermetic container using this manifest, it is as simple as:

\begin{lstlisting}[style=guixcli]
guix shell --container --manifest=manifest.scm -- make test
\end{lstlisting}

\section{Package Transformations: Rewriting the Dependency Graph}

What happens when you need to test your application against a custom Git branch of an upstream dependency, or compile your library with Clang instead of GCC, or replace OpenSSL with a patched security release?

In traditional package managers, this requires editing upstream spec files, maintaining custom binary repositories, or manually rebuilding dozens of packages.

Guix solves this with \textbf{Package Transformations}---dynamic, on-the-fly graph rewrites:

\begin{lstlisting}[style=guixcli]
# 1. Swap source code with a local tarball or directory
guix shell --with-source=python-numpy=./custom-numpy.tar.gz python-numpy

# 2. Pin a dependency to a bleeding-edge Git branch
guix shell --with-branch=libpng=feature-neon inkscape

# 3. Swap one dependency for another across the ENTIRE graph!
guix shell --with-input=openssl=openssl@1.1 node

# 4. Build a package with Clang instead of default GCC
guix shell --with-c-toolchain=ffmpeg=clang-toolchain ffmpeg

# 5. Enable full debug symbols for gdb debugging
guix shell --with-debug-info=glibc --with-debug-info=python python
\end{lstlisting}

\begin{alchemybox}[How Graph Rewriting Works]
When you pass \guixcmd{--with-input=A=B}, Guix performs a recursive transformation on the derivation graph:
\begin{enumerate}
  \item It traverses all packages in the requested dependency DAG.
  \item Every node that lists \texttt{A} as an input has its dependency pointer updated to \texttt{B}.
  \item All downstream packages are automatically rebuilt with their newly derived cryptographic store hashes.
\end{enumerate}
This happens deterministically without mutating any files on your host system.
\end{alchemybox}

\section{Manifest Transformations in Scheme}

You can also embed package transformations directly inside your \texttt{manifest.scm} file so that your entire team inherits them automatically:

\begin{lstlisting}[style=guixscheme]
;; manifest.scm with embedded transformation
(use-modules (guix transformations)
             (gnu packages python)
             (gnu packages python-science))

(define transform-numpy
  (options->transformation
    '((with-branch . "python-numpy=main"))))

(packages->manifest
  (list (transform-numpy python-numpy)
        python-scipy))
\end{lstlisting}

\section{Seamless IDE Integration with \texttt{direnv}}

Manually typing \guixcmd{guix shell} every time you \texttt{cd} into a project directory is good, but automatic activation upon entering the directory is even better.

Using \texttt{direnv}, you can automatically activate your Guix development environment whenever you enter the project directory:

\begin{lstlisting}[style=guixcli]
# In your project root, create a .envrc file:
cat << 'EOF' > .envrc
eval "$(guix shell --search-paths)"
EOF

# Allow direnv to execute in this directory:
direnv allow
\end{lstlisting}

Now, whenever you open VS Code, Emacs, Neovim, or \texttt{cd} into the folder in your terminal, all compiler tools, language servers (like \texttt{pyright}, \texttt{rust-analyzer}, or \texttt{clangd}), and libraries are immediately available in your environment!
